\documentclass[11pt,a4paper]{article}
\usepackage[utf8]{inputenc}
\usepackage[T1]{fontenc}
\usepackage[spanish,es-minimal,es-noshorthands]{babel}
\usepackage{lmodern}
\usepackage[sfdefault,lining]{FiraSans}
\usepackage[a4paper,top=2.3cm,bottom=2.2cm,left=1.8cm,right=1.8cm,headheight=1.3cm,headsep=0.4cm,footskip=1cm]{geometry}
\usepackage{amsmath,amssymb}
\allowdisplaybreaks
\usepackage{icomma}
\usepackage{graphicx}
\usepackage[table]{xcolor}
\usepackage{tikz}
\usepackage{fontawesome5}
\usepackage{ccicons}
\usepackage{fancyhdr}
\usepackage{lastpage}
\usepackage{tasks}
\usepackage{needspace}
\usepackage[most]{tcolorbox}
\usepackage{array}
\usepackage[hidelinks]{hyperref}
\AtBeginDocument{\hypersetup{pdftitle={Ficha de ejercicios de la unidad 5: Límites y continuidad. 2º Bach.},pdfsubject={Matemáticas II 2º Bach.},pdfkeywords={límites, indeterminaciones, infinitésimos equivalentes, límites trigonométricos, parámetros, continuidad, funciones a trozos, teorema de Bolzano, asíntotas, PAU}}}

\definecolor{navy}{HTML}{1B2748}
\definecolor{accent}{HTML}{E8603A}
\definecolor{teal}{HTML}{1F9E90}
\definecolor{soft}{HTML}{EEF1F8}
\definecolor{muted}{HTML}{5D6685}

\setlength{\parindent}{0pt}
\setlength{\parskip}{3pt}
\settasks{label=\alph*), label-format=\bfseries\color{navy}, label-width=1.3em, item-indent=1.8em, column-sep=0.6em, after-item-skip=5pt, before-skip=3pt, after-skip=4pt}

% ---------- capçalera i peu ----------
\newcommand{\llicencia}{\href{https://creativecommons.org/licenses/by-nc-sa/4.0/deed.es}{\ccbyncsa}\ \ Licencia Creative Commons Atribución-NoComercial-CompartirIgual 4.0 Internacional (CC BY-NC-SA 4.0)}
\pagestyle{fancy}
\fancyhf{}
\fancyhead[L]{\raisebox{-0.35\height}{\includegraphics[height=0.95cm]{logo-gva.png}}}
\fancyhead[C]{\small\color{muted}Ficha de ejercicios · Unidad 5: Límites y continuidad · 2\textsuperscript{o} Bach.}
\fancyhead[R]{\raisebox{-0.35\height}{\includegraphics[height=0.95cm]{logo-vermellar.png}}}
\fancyfoot[L]{\scriptsize\color{muted}\llicencia}
\fancyfoot[R]{\small\color{muted}\thepage}
\renewcommand{\headrulewidth}{0.4pt}
\renewcommand{\headrule}{\color{muted!40}\hrule height \headrulewidth}
\fancypagestyle{primera}{\fancyhf{}\renewcommand{\headrulewidth}{0pt}
  \fancyfoot[L]{\scriptsize\color{muted}\llicencia}
  \fancyfoot[R]{\small\color{muted}\thepage}}

% ---------- nivells i exercicis ----------
\newcommand{\estrelles}[1]{{\color{accent}\ifcase#1\or\faStar\faStar[regular]\faStar[regular]\or\faStar\faStar\faStar[regular]\or\faStar\faStar\faStar\fi}}
\newcounter{exer}
\newcommand{\ex}[1]{\par\needspace{4\baselineskip}\medskip\refstepcounter{exer}\noindent{\color{navy}\bfseries\theexer.}\ \ }

\newcommand{\seccio}[2]{\par\needspace{10\baselineskip}\bigskip
  \begin{tcolorbox}[colback=navy,colframe=navy,coltext=white,arc=2mm,boxrule=0pt,left=3mm,right=3mm,top=1.2mm,bottom=1.2mm]
  {\large\bfseries #1\quad #2}
  \end{tcolorbox}\vspace{-2pt}}


\setlength{\emergencystretch}{3em}

% quadre de recordatori (només a la fitxa d'exercicis)
\newcommand{\recorda}[1]{\par\smallskip\begin{tcolorbox}[colback=accent!7,colframe=accent!55,arc=2mm,boxrule=0.5pt,left=2.5mm,right=2.5mm,top=1mm,bottom=1mm]\small{\color{accent}\faLightbulb[regular]}\ \textbf{Recuerda:}\ #1\end{tcolorbox}}
% etiqueta dels exercicis de la PAU d'altres comunitats
\newcommand{\pau}[1]{{\color{teal}\small\bfseries\faUniversity\ PAU · #1.}\ }
\newcommand{\res}[1]{#1}\newcommand{\rest}[1]{#1}

% =====================================================================
% Capçalera FentMates de les fitxes: un color per a cada curs.
% Els colors són els de la pàgina inicial (LEVELS a assets/js/materials.js);
% el logo ve de _eines/logo/pdf/ (python3 _eines/logo/genera.py).
%
% Ús, al final del preàmbul de la unitat:
%   % =====================================================================
% Capçalera FentMates de les fitxes: un color per a cada curs.
% Els colors són els de la pàgina inicial (LEVELS a assets/js/materials.js);
% el logo ve de _eines/logo/pdf/ (python3 _eines/logo/genera.py).
%
% Ús, al final del preàmbul de la unitat:
%   % =====================================================================
% Capçalera FentMates de les fitxes: un color per a cada curs.
% Els colors són els de la pàgina inicial (LEVELS a assets/js/materials.js);
% el logo ve de _eines/logo/pdf/ (python3 _eines/logo/genera.py).
%
% Ús, al final del preàmbul de la unitat:
%   \input{../comu/fentmates.tex}
%   \fmcurs{1eso}
% i, al començament del document:
%   \capcalera{Matemàtiques}{1r ESO}{Fitxa d'exercicis · Unitat 2}{Divisibilitat}
% Opcional abans de \input: \def\fmneed{14} (espai mínim abans de cada secció).
% =====================================================================
\makeatletter
\def\fm@logodir{../../logo/pdf}
\providecommand{\fmneed}{10}
\def\fm@cursnom{}

% clau del curs, color i color del text damunt del color
\newcommand{\fm@curs}[3]{\@namedef{fm@color@#1}{#2}\@namedef{fm@text@#1}{#3}}
\fm@curs{1eso}{E6B84A}{navy}
\fm@curs{2eso}{3B82F6}{white}
\fm@curs{3eso}{FF7A50}{white}
\fm@curs{4eso-a}{A78BFA}{white}
\fm@curs{4eso-b}{6D28D9}{white}
\fm@curs{1batx-mi}{9B1C1C}{white}
\fm@curs{1batx-ccss}{EF6B6B}{white}
\fm@curs{2batx-mii}{1E6B47}{white}
\fm@curs{2batx-ccss}{2BB3A3}{white}

\newcommand{\fmcurs}[1]{%
  \@ifundefined{fm@color@#1}{\PackageError{fentmates}{Curs desconegut: #1}{}}{}%
  \edef\fm@tmp{\noexpand\definecolor{curs}{HTML}{\@nameuse{fm@color@#1}}%
               \noexpand\colorlet{curstext}{\@nameuse{fm@text@#1}}}\fm@tmp
  \colorlet{cursfosc}{curs!72!black}%
  \def\fm@logo{\fm@logodir/vertical-#1.pdf}%
  \def\fm@marca{\fm@logodir/marca-#1.pdf}}

% ---------- capçalera de la primera pàgina ----------
% #1 matèria (petit), #2 curs, #3 tipus de document i unitat, #4 títol
\newcommand{\capcalera}[4]{%
  \gdef\fm@cursnom{#2}%
  \par\noindent
  \begin{tcolorbox}[colback=curs,colframe=curs,boxrule=0pt,arc=2.2mm,width=3.5cm,height=2.15cm,
      valign=center,left=3mm,right=2mm,top=1mm,bottom=1mm,nobeforeafter,box align=center]
    \color{curstext}{\footnotesize #1}\par\vspace{1pt}{\fontsize{19}{21}\selectfont\bfseries #2\par}%
  \end{tcolorbox}\hspace{5mm}%
  \begin{minipage}[c]{\dimexpr\linewidth-3.5cm-5mm-2.9cm\relax}
    {\color{cursfosc}\small\bfseries\MakeUppercase{#3}}\par\vspace{2pt}
    {\color{navy}\fontsize{20}{23}\selectfont\bfseries #4\par}
  \end{minipage}\hfill
  \begin{minipage}[c]{2.6cm}\raggedleft\includegraphics[width=2.6cm]{\fm@logo}\end{minipage}\par
  \vspace{2pt}}

% ---------- barres de secció del color del curs ----------
\renewcommand{\seccio}[2]{\par\needspace{\fmneed\baselineskip}\bigskip
  \begin{tcolorbox}[colback=curs,colframe=curs,coltext=curstext,arc=2mm,boxrule=0pt,left=3mm,right=3mm,top=1.2mm,bottom=1.2mm]
  {\large\bfseries #1\quad #2}
  \end{tcolorbox}\vspace{-2pt}}

% ---------- estrelles de dificultat del color del curs (a Batxillerat no se'n posen) ----------
\@ifundefined{estrelles}{}{%
  \renewcommand{\estrelles}[1]{{\color{curs}\ifcase#1\or\faStar\faStar[regular]\faStar[regular]\or\faStar\faStar\faStar[regular]\or\faStar\faStar\faStar\fi}}}

% ---------- capçalera i peu de les pàgines ----------
% primera pàgina: logos institucionals al peu
\fancypagestyle{primera}{\fancyhf{}\renewcommand{\headrulewidth}{0pt}
  \fancyfoot[L]{\raisebox{-0.32\height}{\includegraphics[height=0.78cm]{logo-gva.png}}\hspace{5mm}%
    \raisebox{-0.32\height}{\includegraphics[height=0.78cm]{logo-vermellar.png}}}
  \fancyfoot[R]{\scriptsize\color{muted}\href{https://creativecommons.org/licenses/by-nc-sa/4.0/deed.ca}{\ccbyncsa}\ \ CC BY-NC-SA 4.0 · fentmates.org\qquad\small\thepage}}
% la resta: el símbol a l'esquerra i el curs a la dreta, amb el títol al centre com abans
\fancyhead[L]{\raisebox{-0.3\height}{\includegraphics[height=0.62cm]{\fm@marca}}}
\fancyhead[R]{\small\bfseries\color{cursfosc}\fm@cursnom}
\renewcommand{\headrule}{\color{curs!45}\hrule height \headrulewidth}
\makeatother

%   \fmcurs{1eso}
% i, al començament del document:
%   \capcalera{Matemàtiques}{1r ESO}{Fitxa d'exercicis · Unitat 2}{Divisibilitat}
% Opcional abans de \input: \def\fmneed{14} (espai mínim abans de cada secció).
% =====================================================================
\makeatletter
\def\fm@logodir{../../logo/pdf}
\providecommand{\fmneed}{10}
\def\fm@cursnom{}

% clau del curs, color i color del text damunt del color
\newcommand{\fm@curs}[3]{\@namedef{fm@color@#1}{#2}\@namedef{fm@text@#1}{#3}}
\fm@curs{1eso}{E6B84A}{navy}
\fm@curs{2eso}{3B82F6}{white}
\fm@curs{3eso}{FF7A50}{white}
\fm@curs{4eso-a}{A78BFA}{white}
\fm@curs{4eso-b}{6D28D9}{white}
\fm@curs{1batx-mi}{9B1C1C}{white}
\fm@curs{1batx-ccss}{EF6B6B}{white}
\fm@curs{2batx-mii}{1E6B47}{white}
\fm@curs{2batx-ccss}{2BB3A3}{white}

\newcommand{\fmcurs}[1]{%
  \@ifundefined{fm@color@#1}{\PackageError{fentmates}{Curs desconegut: #1}{}}{}%
  \edef\fm@tmp{\noexpand\definecolor{curs}{HTML}{\@nameuse{fm@color@#1}}%
               \noexpand\colorlet{curstext}{\@nameuse{fm@text@#1}}}\fm@tmp
  \colorlet{cursfosc}{curs!72!black}%
  \def\fm@logo{\fm@logodir/vertical-#1.pdf}%
  \def\fm@marca{\fm@logodir/marca-#1.pdf}}

% ---------- capçalera de la primera pàgina ----------
% #1 matèria (petit), #2 curs, #3 tipus de document i unitat, #4 títol
\newcommand{\capcalera}[4]{%
  \gdef\fm@cursnom{#2}%
  \par\noindent
  \begin{tcolorbox}[colback=curs,colframe=curs,boxrule=0pt,arc=2.2mm,width=3.5cm,height=2.15cm,
      valign=center,left=3mm,right=2mm,top=1mm,bottom=1mm,nobeforeafter,box align=center]
    \color{curstext}{\footnotesize #1}\par\vspace{1pt}{\fontsize{19}{21}\selectfont\bfseries #2\par}%
  \end{tcolorbox}\hspace{5mm}%
  \begin{minipage}[c]{\dimexpr\linewidth-3.5cm-5mm-2.9cm\relax}
    {\color{cursfosc}\small\bfseries\MakeUppercase{#3}}\par\vspace{2pt}
    {\color{navy}\fontsize{20}{23}\selectfont\bfseries #4\par}
  \end{minipage}\hfill
  \begin{minipage}[c]{2.6cm}\raggedleft\includegraphics[width=2.6cm]{\fm@logo}\end{minipage}\par
  \vspace{2pt}}

% ---------- barres de secció del color del curs ----------
\renewcommand{\seccio}[2]{\par\needspace{\fmneed\baselineskip}\bigskip
  \begin{tcolorbox}[colback=curs,colframe=curs,coltext=curstext,arc=2mm,boxrule=0pt,left=3mm,right=3mm,top=1.2mm,bottom=1.2mm]
  {\large\bfseries #1\quad #2}
  \end{tcolorbox}\vspace{-2pt}}

% ---------- estrelles de dificultat del color del curs (a Batxillerat no se'n posen) ----------
\@ifundefined{estrelles}{}{%
  \renewcommand{\estrelles}[1]{{\color{curs}\ifcase#1\or\faStar\faStar[regular]\faStar[regular]\or\faStar\faStar\faStar[regular]\or\faStar\faStar\faStar\fi}}}

% ---------- capçalera i peu de les pàgines ----------
% primera pàgina: logos institucionals al peu
\fancypagestyle{primera}{\fancyhf{}\renewcommand{\headrulewidth}{0pt}
  \fancyfoot[L]{\raisebox{-0.32\height}{\includegraphics[height=0.78cm]{logo-gva.png}}\hspace{5mm}%
    \raisebox{-0.32\height}{\includegraphics[height=0.78cm]{logo-vermellar.png}}}
  \fancyfoot[R]{\scriptsize\color{muted}\href{https://creativecommons.org/licenses/by-nc-sa/4.0/deed.ca}{\ccbyncsa}\ \ CC BY-NC-SA 4.0 · fentmates.org\qquad\small\thepage}}
% la resta: el símbol a l'esquerra i el curs a la dreta, amb el títol al centre com abans
\fancyhead[L]{\raisebox{-0.3\height}{\includegraphics[height=0.62cm]{\fm@marca}}}
\fancyhead[R]{\small\bfseries\color{cursfosc}\fm@cursnom}
\renewcommand{\headrule}{\color{curs!45}\hrule height \headrulewidth}
\makeatother

%   \fmcurs{1eso}
% i, al començament del document:
%   \capcalera{Matemàtiques}{1r ESO}{Fitxa d'exercicis · Unitat 2}{Divisibilitat}
% Opcional abans de \input: \def\fmneed{14} (espai mínim abans de cada secció).
% =====================================================================
\makeatletter
\def\fm@logodir{../../logo/pdf}
\providecommand{\fmneed}{10}
\def\fm@cursnom{}

% clau del curs, color i color del text damunt del color
\newcommand{\fm@curs}[3]{\@namedef{fm@color@#1}{#2}\@namedef{fm@text@#1}{#3}}
\fm@curs{1eso}{E6B84A}{navy}
\fm@curs{2eso}{3B82F6}{white}
\fm@curs{3eso}{FF7A50}{white}
\fm@curs{4eso-a}{A78BFA}{white}
\fm@curs{4eso-b}{6D28D9}{white}
\fm@curs{1batx-mi}{9B1C1C}{white}
\fm@curs{1batx-ccss}{EF6B6B}{white}
\fm@curs{2batx-mii}{1E6B47}{white}
\fm@curs{2batx-ccss}{2BB3A3}{white}

\newcommand{\fmcurs}[1]{%
  \@ifundefined{fm@color@#1}{\PackageError{fentmates}{Curs desconegut: #1}{}}{}%
  \edef\fm@tmp{\noexpand\definecolor{curs}{HTML}{\@nameuse{fm@color@#1}}%
               \noexpand\colorlet{curstext}{\@nameuse{fm@text@#1}}}\fm@tmp
  \colorlet{cursfosc}{curs!72!black}%
  \def\fm@logo{\fm@logodir/vertical-#1.pdf}%
  \def\fm@marca{\fm@logodir/marca-#1.pdf}}

% ---------- capçalera de la primera pàgina ----------
% #1 matèria (petit), #2 curs, #3 tipus de document i unitat, #4 títol
\newcommand{\capcalera}[4]{%
  \gdef\fm@cursnom{#2}%
  \par\noindent
  \begin{tcolorbox}[colback=curs,colframe=curs,boxrule=0pt,arc=2.2mm,width=3.5cm,height=2.15cm,
      valign=center,left=3mm,right=2mm,top=1mm,bottom=1mm,nobeforeafter,box align=center]
    \color{curstext}{\footnotesize #1}\par\vspace{1pt}{\fontsize{19}{21}\selectfont\bfseries #2\par}%
  \end{tcolorbox}\hspace{5mm}%
  \begin{minipage}[c]{\dimexpr\linewidth-3.5cm-5mm-2.9cm\relax}
    {\color{cursfosc}\small\bfseries\MakeUppercase{#3}}\par\vspace{2pt}
    {\color{navy}\fontsize{20}{23}\selectfont\bfseries #4\par}
  \end{minipage}\hfill
  \begin{minipage}[c]{2.6cm}\raggedleft\includegraphics[width=2.6cm]{\fm@logo}\end{minipage}\par
  \vspace{2pt}}

% ---------- barres de secció del color del curs ----------
\renewcommand{\seccio}[2]{\par\needspace{\fmneed\baselineskip}\bigskip
  \begin{tcolorbox}[colback=curs,colframe=curs,coltext=curstext,arc=2mm,boxrule=0pt,left=3mm,right=3mm,top=1.2mm,bottom=1.2mm]
  {\large\bfseries #1\quad #2}
  \end{tcolorbox}\vspace{-2pt}}

% ---------- estrelles de dificultat del color del curs (a Batxillerat no se'n posen) ----------
\@ifundefined{estrelles}{}{%
  \renewcommand{\estrelles}[1]{{\color{curs}\ifcase#1\or\faStar\faStar[regular]\faStar[regular]\or\faStar\faStar\faStar[regular]\or\faStar\faStar\faStar\fi}}}

% ---------- capçalera i peu de les pàgines ----------
% primera pàgina: logos institucionals al peu
\fancypagestyle{primera}{\fancyhf{}\renewcommand{\headrulewidth}{0pt}
  \fancyfoot[L]{\raisebox{-0.32\height}{\includegraphics[height=0.78cm]{logo-gva.png}}\hspace{5mm}%
    \raisebox{-0.32\height}{\includegraphics[height=0.78cm]{logo-vermellar.png}}}
  \fancyfoot[R]{\scriptsize\color{muted}\href{https://creativecommons.org/licenses/by-nc-sa/4.0/deed.ca}{\ccbyncsa}\ \ CC BY-NC-SA 4.0 · fentmates.org\qquad\small\thepage}}
% la resta: el símbol a l'esquerra i el curs a la dreta, amb el títol al centre com abans
\fancyhead[L]{\raisebox{-0.3\height}{\includegraphics[height=0.62cm]{\fm@marca}}}
\fancyhead[R]{\small\bfseries\color{cursfosc}\fm@cursnom}
\renewcommand{\headrule}{\color{curs!45}\hrule height \headrulewidth}
\makeatother

\fmcurs{2batx-mii}

\begin{document}
\thispagestyle{primera}

\capcalera{Matemáticas II}{2\textsuperscript{o} Bach.}{Ficha de ejercicios · Unidad 5}{\raggedright Límites y continuidad}
\vspace{10pt}
Nombre y apellidos: \rule{7.2cm}{0.4pt}\hfill Grupo: \rule{1.6cm}{0.4pt}\hfill Fecha: \rule{2.2cm}{0.4pt}
\vspace{8pt}

\begin{tcolorbox}[colback=soft,colframe=soft,arc=2mm,left=3mm,right=3mm,top=1.5mm,bottom=1.5mm]
\small Resuelve los ejercicios en el cuaderno justificando cada paso. Ningún límite de esta ficha necesita derivadas (la regla de L'Hôpital la veremos más adelante). Los ejercicios marcados con {\color{teal}\bfseries\faUniversity\ PAU} son de las pruebas de acceso a la universidad de otras comunidades autónomas, traducidos al castellano (solo los apartados que corresponden a esta unidad).
\end{tcolorbox}

% =====================================================
\seccio{1}{Límites en el infinito}

\ex{2} Calcula los límites siguientes:\begin{tasks}(3)
  \task $\displaystyle\lim_{x\to -\infty}(2x^5-x^2)$
  \task $\displaystyle\lim_{x\to +\infty}\dfrac{4x^3-x}{2x^3+7}$
  \task $\displaystyle\lim_{x\to -\infty}\dfrac{1-3x^2}{x+2}$
  \task $\displaystyle\lim_{x\to +\infty}\dfrac{x^2+1}{x^3-2}$
  \task $\displaystyle\lim_{x\to +\infty}\dfrac{2^{x+1}+5}{2^{x}-3}$
  \task $\displaystyle\lim_{x\to -\infty}\dfrac{3e^{2x}-e^{x}}{e^{2x}+4}$
\end{tasks}\recorda{Cuando $x\to\pm\infty$, en un cociente de polinomios mandan los términos de mayor grado. Con exponenciales, dividimos por la exponencial que crece más deprisa (y recuerda que $e^{x}\to0$ cuando $x\to-\infty$).}

\ex{2} Calcula los límites siguientes con raíces. Cuidado con el signo cuando $x\to-\infty$:\begin{tasks}(3)
  \task $\displaystyle\lim_{x\to +\infty}\dfrac{\sqrt{9x^2+1}}{2x-5}$
  \task $\displaystyle\lim_{x\to -\infty}\dfrac{\sqrt{9x^2+1}}{2x-5}$
  \task $\displaystyle\lim_{x\to -\infty}\dfrac{3x+\sqrt{x^2+5}}{2x+1}$
  \task $\displaystyle\lim_{x\to -\infty}\dfrac{\sqrt[3]{8x^3-1}}{x}$
  \task $\displaystyle\lim_{x\to +\infty}\dfrac{\sqrt{x^3+x}}{x+4}$
  \task $\displaystyle\lim_{x\to -\infty}\dfrac{\sqrt{x^2-5x+3}}{3x-2}$
\end{tasks}\recorda{$\sqrt{x^2}=|x|$: cuando $x\to-\infty$, $\sqrt{a x^2+\dots}\approx-\sqrt a\,x$. Si dudas, haz el cambio $x=-t$ con $t\to+\infty$.}

\ex{2} Ordena de menor a mayor orden de infinito (cuando $x\to+\infty$) \ $x^3,\ 2^x,\ \ln x,\ \sqrt x,\ e^x,\ x^{10},\ (\ln x)^3$, \ y calcula:\begin{tasks}(3)
  \task $\displaystyle\lim_{x\to +\infty}\dfrac{x^3}{e^x}$
  \task $\displaystyle\lim_{x\to +\infty}\dfrac{\ln x}{x^2}$
  \task $\displaystyle\lim_{x\to +\infty}\dfrac{2^x+x^4}{3\cdot2^x}$
  \task $\displaystyle\lim_{x\to +\infty}\left(x^{10}-1{,}5^x\right)$
  \task $\displaystyle\lim_{x\to +\infty}\dfrac{3^x+2^x}{3^x-x^2}$
  \task $\displaystyle\lim_{x\to +\infty}\dfrac{(\ln x)^3}{\sqrt x}$
\end{tasks}

\ex{2} Calcula, utilizando las propiedades de los logaritmos:\begin{tasks}(2)
  \task $\displaystyle\lim_{x\to +\infty}\big(\ln(2x+1)-\ln(x+3)\big)$
  \task $\displaystyle\lim_{x\to +\infty}\big(\ln(x^2+1)-\ln x\big)$
\end{tasks}

\ex{2} \pau{Andalucía, ordinaria 2023} Considera la función $f(x)=\dfrac{1}{e^x+e^{-x}}$. Calcula $\displaystyle\lim_{x\to+\infty}x^2f(x)$.


% =====================================================
\seccio{2}{Indeterminaciones en el infinito}

\ex{2} Indeterminación $\infty-\infty$:\begin{tasks}(3)
  \task $\displaystyle\lim_{x\to +\infty}\left(\dfrac{x^2}{x+1}-x\right)$
  \task $\displaystyle\lim_{x\to +\infty}\left(\dfrac{x^3}{x^2+1}-\dfrac{x^2}{x+1}\right)$
  \task $\displaystyle\lim_{x\to +\infty}\left(\sqrt{x^2+8x}-x\right)$
  \task $\displaystyle\lim_{x\to +\infty}\left(\sqrt{x^2+3x}-\sqrt{x^2-x}\right)$
  \task $\displaystyle\lim_{x\to -\infty}\left(\sqrt{x^2+2x}+x\right)$
  \task $\displaystyle\lim_{x\to +\infty}\left(\sqrt{x^2+1}-3x\right)$
\end{tasks}\recorda{Con fracciones, hacemos la resta. Con raíces, multiplicamos y dividimos por el \textbf{conjugado}. Antes, comprueba si los términos principales se cancelan: si no lo hacen, no hay indeterminación real.}

\ex{3} Calcula $\displaystyle\lim_{x\to+\infty}\left(\sqrt[3]{x^3+3x^2}-x\right)$. \ \textit{Indicación}: $A^3-B^3=(A-B)(A^2+AB+B^2)$.

\ex{2} Indeterminación $0\cdot\infty$:\begin{tasks}(2)
  \task $\displaystyle\lim_{x\to +\infty}x\left(\sqrt{x^2+4}-x\right)$
  \task $\displaystyle\lim_{x\to +\infty}(2x+1)\left(\sqrt{x+1}-\sqrt x\right)$
  \task $\displaystyle\lim_{x\to +\infty}x\ln\left(1+\dfrac2x\right)$
  \task $\displaystyle\lim_{x\to +\infty}x^3e^{-x}$
\end{tasks}\recorda{Transformamos el producto en un cociente ($f\cdot g=\frac{f}{1/g}$) u operamos. Un exponente se puede poner dentro del logaritmo: $x\ln A=\ln A^{x}$.}

\ex{2} Indeterminación $1^{\infty}$:\begin{tasks}(3)
  \task $\displaystyle\lim_{x\to +\infty}\left(\dfrac{x+2}{x-1}\right)^{2x}$
  \task $\displaystyle\lim_{x\to +\infty}\left(\dfrac{2x-1}{2x+3}\right)^{x+1}$
  \task $\displaystyle\lim_{x\to +\infty}\left(\dfrac{x^2+3}{x^2-1}\right)^{x^2}$
  \task $\displaystyle\lim_{x\to +\infty}\left(\dfrac{3x+1}{x+2}\right)^{x}$
  \task $\displaystyle\lim_{x\to -\infty}\left(\dfrac{x-4}{x}\right)^{3x}$
  \task $\displaystyle\lim_{x\to +\infty}\left(\dfrac{x^2-x+1}{x^2+2}\right)^{\frac{x^2+1}{x-3}}$
\end{tasks}\recorda{Si $f\to1$ y $g\to\pm\infty$: \ $\lim f^{\,g}=e^{\lim\,(f-1)\,g}$. Si la base no tiende a 1, no hay indeterminación.}

\ex{3} Indeterminaciones $\infty^0$. Calcula tomando logaritmos:\begin{tasks}(2)
  \task $\displaystyle\lim_{x\to +\infty}x^{1/x}$
  \task $\displaystyle\lim_{x\to +\infty}(3x)^{\frac{2}{\ln x}}$
\end{tasks}

\ex{2} \pau{Aragón, ordinaria 2024} Calcula justificadamente $\displaystyle\lim_{x\to+\infty}\left[\sqrt{x^2+5}-(x+2)\right]$.

\ex{2} \pau{Aragón, ordinaria 2022} Calcula $\displaystyle\lim_{x\to+\infty}\left[\sqrt{3x^2-2}-\left(\sqrt3\,x+5\right)\right]$.

\ex{2} \pau{Navarra, extraordinaria 2021} Calcula $\displaystyle\lim_{x\to+\infty}\dfrac{\sqrt x}{\sqrt{3x^3+2x^2}-\sqrt{3x^3}}$.

\ex{2} \pau{Murcia, ordinaria 2024} Calcula:\begin{tasks}(2)
  \task $\displaystyle\lim_{x\to +\infty}\left(\sqrt{x+9}-\sqrt{x-9}\right)$
  \task $\displaystyle\lim_{x\to +\infty}\dfrac{\ln x}{\sqrt x}$
\end{tasks}

\ex{2} \pau{Castilla-La Mancha, ordinaria 2023} Calcula $\displaystyle\lim_{x\to+\infty}\left(\dfrac{5x+1}{5x}\right)^{x^2}$.

\ex{2} \pau{Navarra, extraordinaria 2022} Calcula $\displaystyle\lim_{x\to+\infty}\left(\dfrac{x^2}{x^2+x-1}\right)^{2x-1}$.

\ex{2} \pau{La Rioja, extraordinaria 2023} Calcula $\displaystyle\lim_{x\to+\infty}(1+x)^{1/x}$.

\ex{2} \pau{Islas Baleares, extraordinaria 2021} El número de individuos, en millones, de una población es $P(t)=\dfrac{15+t^2}{(t+1)^2}$, donde $t\geq0$ son los años transcurridos desde el 1 de enero de 2000.\begin{tasks}(1)
  \task ¿Qué población había el 1 de enero de 2000?
  \task Calcula el tamaño de la población a largo plazo.
\end{tasks}


% =====================================================
\seccio{3}{Límites en un número e indeterminaciones}

\ex{2} Calcula los límites laterales y di si existe el límite:\begin{tasks}(2)
  \task $\displaystyle\lim_{x\to 2}\dfrac{x+1}{(x-2)^2}$
  \task $\displaystyle\lim_{x\to 3}\dfrac{x-3}{x^3-5x^2+3x+9}$
\end{tasks}

\ex{2} Indeterminación $\frac00$ con polinomios:\begin{tasks}(3)
  \task $\displaystyle\lim_{x\to 2}\dfrac{x^3-8}{x-2}$
  \task $\displaystyle\lim_{x\to 1}\dfrac{x^4-1}{x^2-x}$
  \task $\displaystyle\lim_{x\to 2}\dfrac{x^4-4x^3+16x-16}{x^3-6x^2+12x-8}$
\end{tasks}\recorda{Si $P(a)=Q(a)=0$, ambos son divisibles por $(x-a)$: factorizamos (Ruffini), simplificamos y volvemos a sustituir.}

\ex{2} Indeterminación $\frac00$ con raíces:\begin{tasks}(3)
  \task $\displaystyle\lim_{x\to 4}\dfrac{\sqrt{x+5}-3}{x-4}$
  \task $\displaystyle\lim_{x\to 3}\dfrac{\sqrt{x+1}-\sqrt{2x-2}}{x^2-9}$
  \task $\displaystyle\lim_{x\to 0}\dfrac{\sqrt{1+x}-1}{\sqrt{4+x}-2}$
\end{tasks}

\ex{2} Indeterminaciones $\infty-\infty$ y $1^\infty$ en un número:\begin{tasks}(3)
  \task $\displaystyle\lim_{x\to 1}\left(\dfrac{1}{x-1}-\dfrac{2}{x^2-1}\right)$
  \task $\displaystyle\lim_{x\to 2}\left(\dfrac{x}{x-2}-\dfrac{8}{x^2-4}\right)$
  \task $\displaystyle\lim_{x\to 0}(1+3x)^{2/x}$
  \task $\displaystyle\lim_{x\to 2}\left(\dfrac{x+2}{2x}\right)^{\frac{1}{x-2}}$
  \task $\displaystyle\lim_{x\to 1}x^{\frac{3}{x-1}}$
  \task $\displaystyle\lim_{x\to 1}(2x-1)^{\frac{1}{x-1}}$
\end{tasks}

\ex{2} \pau{Castilla-La Mancha, extraordinaria 2024} Calcula $\displaystyle\lim_{x\to-2}\dfrac{x^3+2x^2+x+2}{x^2+2x}$.

\ex{2} \pau{Castilla-La Mancha, extraordinaria 2023} Calcula $\displaystyle\lim_{x\to3}\dfrac{x^3-3x^2+3x-9}{3x-9}$.

\ex{2} \pau{Castilla y León, ordinaria 2023} Calcula $\displaystyle\lim_{x\to2}\dfrac{\sqrt{x^3+x-1}-\sqrt{x^3+1}}{x-2}$.

\ex{2} \pau{Murcia, ordinaria 2023} Calcula $\displaystyle\lim_{x\to0}\dfrac{\sqrt{9+x}-\sqrt{9-x}}{3x}$.


% =====================================================
\seccio{4}{Límites trigonométricos e infinitésimos equivalentes}

\ex{2} Calcula, utilizando infinitésimos equivalentes:\begin{tasks}(3)
  \task $\displaystyle\lim_{x\to 0}\dfrac{\sin 4x}{3x}$
  \task $\displaystyle\lim_{x\to 0}\dfrac{x\tan 2x}{1-\cos x}$
  \task $\displaystyle\lim_{x\to 0}\dfrac{\ln(1+x^2)}{\sin^2 3x}$
  \task $\displaystyle\lim_{x\to 0}\dfrac{\sin 3x}{\tan 5x}$
  \task $\displaystyle\lim_{x\to 0}\dfrac{1-\cos 2x}{x\sin x}$
  \task $\displaystyle\lim_{x\to 0}\dfrac{e^{2x}-1}{\ln(1+3x)}$
\end{tasks}\recorda{Si $u\to0$: \ $\sin u\sim u$, \ $\tan u\sim u$, \ $1-\cos u\sim\frac{u^2}{2}$, \ $e^u-1\sim u$, \ $\ln(1+u)\sim u$. Solo se pueden sustituir en \textbf{productos y cocientes}, nunca en sumas o restas.}

\ex{2} Calcula (si el punto no es 0, haz primero un cambio de variable):\begin{tasks}(3)
  \task $\displaystyle\lim_{x\to 0}\dfrac{e^{x}-e^{-x}}{\sin x}$
  \task $\displaystyle\lim_{x\to \frac\pi2}\dfrac{\cos x}{2x-\pi}$
  \task $\displaystyle\lim_{x\to \pi}\dfrac{\sin x}{x-\pi}$
  \task $\displaystyle\lim_{x\to 1}\dfrac{\ln x}{x^2-1}$
  \task $\displaystyle\lim_{x\to 0}\dfrac{\sin x-\tan x}{x^2\sin x}$
  \task $\displaystyle\lim_{x\to 0}(\cos x)^{1/x^2}$
\end{tasks}

\ex{2} Calcula, utilizando que el producto de una función acotada por una que tiende a 0 tiende a 0:\begin{tasks}(3)
  \task $\displaystyle\lim_{x\to -\infty}\dfrac{3x+\cos x}{x}$
  \task $\displaystyle\lim_{x\to 0}x^2\cos\dfrac1x$
  \task $\displaystyle\lim_{x\to +\infty}x\sin\dfrac1x$
\end{tasks}

\ex{2} \pau{Murcia, ordinaria 2023} Calcula $\displaystyle\lim_{x\to0}\dfrac{\cos(2x)-1}{x\sin x}$.

\ex{2} \pau{Castilla y León, extraordinaria 2022} Calcula $\displaystyle\lim_{x\to0}\dfrac{\ln(1+x)}{e^x-1}$.

\ex{2} \pau{Galicia, ordinaria 2026} Calcula los límites siguientes:\begin{tasks}(2)
  \task $\displaystyle\lim_{x\to 1}\dfrac{\ln x}{x}$
  \task $\displaystyle\lim_{x\to 1}\dfrac{x-1}{\ln x}$
  \task $\displaystyle\lim_{x\to 0}\dfrac{\tan x}{\sin x}$
  \task $\displaystyle\lim_{x\to 0^+}\sqrt x\,\ln x$
\end{tasks}

\ex{2} \pau{La Rioja, extraordinaria 2023} Calcula $\displaystyle\lim_{x\to0}(\cos2x)^{3/x^2}$.

\ex{3} \pau{Aragón, extraordinaria 2021} Calcula $\displaystyle\lim_{x\to1}\left(\sin\dfrac{\pi x}{2}\right)^{\frac{1}{(1-x)^2}}$.

\ex{2} \pau{Navarra, extraordinaria 2021} Calcula $\displaystyle\lim_{x\to\infty}x\sin\dfrac1x$.

\ex{2} \pau{Murcia, extraordinaria 2025} Calcula $\displaystyle\lim_{x\to+\infty}x\left(e^{1/x}-1\right)$.


% =====================================================
\seccio{5}{Límites con parámetros}

\ex{2} Calcula el valor de los parámetros para que:\begin{tasks}(1)
  \task $\displaystyle\lim_{x\to +\infty}\dfrac{ax^2-3x}{2x^2+5}=4$
  \task $\displaystyle\lim_{x\to +\infty}\dfrac{ax^3+bx^2+1}{2x^2-x}=3$
  \task $\displaystyle\lim_{x\to +\infty}\left(\dfrac{2x+3}{2x+a}\right)^{4x}=e^2$
  \task $\displaystyle\lim_{x\to 0}\dfrac{e^{mx}-1}{\sin 2x}=\dfrac32$
\end{tasks}\recorda{Calculamos el límite como si el parámetro fuera un número e imponemos la condición.}

\ex{2} Halla el valor de $a$ para que el límite sea finito, y calcúlalo en cada caso:\begin{tasks}(2)
  \task $\displaystyle\lim_{x\to +\infty}\left(\sqrt{ax^2+x}-2x\right)$
  \task $\displaystyle\lim_{x\to -3}\dfrac{x^2+ax+3}{x^2+3x}$
\end{tasks}

\ex{2} Halla $a$ y $b$ para que $\displaystyle\lim_{x\to1}\dfrac{x^2+ax+b}{x-1}=3$.

\ex{2} Halla $a$ para que:\begin{tasks}(2)
  \task $\displaystyle\lim_{x\to +\infty}\left(\sqrt{x^2+ax}-\sqrt{x^2-x}\right)=2$
  \task $\displaystyle\lim_{x\to 0}\dfrac{1-\cos(ax)}{x\sin x}=8$
\end{tasks}

\ex{3} Discute, según los valores del parámetro $a$, los límites siguientes:\begin{tasks}(1)
  \task $\displaystyle\lim_{x\to +\infty}\dfrac{(a-1)x^2+x}{3x+1}$
  \task $\displaystyle\lim_{x\to +\infty}\left(\dfrac{x+1}{x+2}\right)^{ax^2+x}$
\end{tasks}

\ex{2} \pau{Andalucía, extraordinaria 2022} Calcula $a$ sabiendo que $\displaystyle\lim_{x\to+\infty}\dfrac{ax}{(\ln x)^2+2x}=1$.

\ex{3} \pau{Aragón, extraordinaria 2025} Calcula (si existe), en función del valor de $k\in\mathbb Z$, el límite \ $\displaystyle\lim_{x\to1}\dfrac{x^4+x^3-x^2-x}{\left(x^2-1\right)^{2k}}$.

\ex{2} \pau{Castilla-La Mancha, extraordinaria 2022} Halla razonadamente $a$ y $b$ para que la función $f(x)=\dfrac{ax+1}{2x+b}$ tenga una discontinuidad de salto infinito en $x=1$ y tienda a 2 cuando $x\to+\infty$.


% =====================================================
\seccio{6}{Continuidad de funciones}

\ex{2} Estudia la continuidad en $x=0$ de las funciones siguientes y clasifica la discontinuidad, si la hay:\begin{tasks}(1)
  \task $f(x)=\begin{cases}\dfrac{\sin 3x}{x}&\text{si }x\neq0\\[3pt]3&\text{si }x=0\end{cases}$
  \task $g(x)=\dfrac{1}{1+e^{1/x}}$
  \task $h(x)=x\cos\dfrac1x$. ¿Qué valor habría que dar a $h(0)$ para que fuera continua?
\end{tasks}\recorda{$f$ es continua en $a$ si $\exists f(a)$, $\exists\lim_{x\to a}f(x)$ y son iguales. Discontinuidades: evitable (el límite existe), de salto finito (laterales finitos y distintos), de salto infinito (un lateral infinito) y asintótica (ambos infinitos).}

\ex{2} Estudia la continuidad de las funciones siguientes:\begin{tasks}(1)
  \task $f(x)=\begin{cases}\cos x&\text{si }x\leq0\\[3pt]\dfrac{e^x-1}{x}&\text{si }0<x<2\\[3pt]\log_2(x+2)&\text{si }x\geq2\end{cases}$
  \task $g(x)=\begin{cases}\dfrac{3}{x+2}&\text{si }x<-1\\[3pt]2-|x|&\text{si }-1\leq x\leq1\\[3pt]\sqrt{x+3}-1&\text{si }x>1\end{cases}$
  \task $h(x)=\begin{cases}\dfrac{\sin 2x}{x}&\text{si }x<0\\[3pt]x+2&\text{si }0\leq x<\pi\\[3pt]\cos x+\pi+3&\text{si }x\geq\pi\end{cases}$
\end{tasks}\recorda{Primero, cada trozo en su intervalo abierto (denominadores, logaritmos, raíces, $\tan$\ldots\ solo cuentan si el punto está dentro del trozo). Después, los puntos donde cambia la expresión.}

\ex{2} Estudia la continuidad de:\begin{tasks}(3)
  \task $f(x)=\dfrac{x}{1+e^{1/x}}$
  \task $g(x)=\dfrac{x^2-4}{|x-2|}$
  \task $h(x)=e^{\frac{1}{x-1}}$
\end{tasks}

\ex{2} \pau{Castilla-La Mancha, extraordinaria 2024} Estudia la continuidad de la función y, si las tiene, clasifica sus discontinuidades:\[f(x)=\begin{cases}x^2+2x+1&\text{si }x<3\\[3pt]\dfrac{2x}{x-4}&\text{si }x\geq3\end{cases}\]

\ex{2} \pau{Castilla-La Mancha, extraordinaria 2020} Estudia la continuidad de $f$ en $x=1$ y en $x=2$ y clasifica en ellos el tipo de discontinuidad, si la hay:\[f(x)=\begin{cases}2^{x-1}&\text{si }x\leq1\\[3pt]x-2&\text{si }1<x<2\\[3pt]\ln(x-1)&\text{si }x\geq2\end{cases}\]

\ex{2} \pau{Asturias, ordinaria 2025} Estudia si la función es continua en su dominio:\[f(x)=\begin{cases}\dfrac x2&\text{si }0\leq x<4\\[3pt]3-(x-5)^2&\text{si }4\leq x\end{cases}\]

\ex{2} \pau{Madrid, ordinaria 2023} Dada la función\[f(x)=\begin{cases}\dfrac{x^2}{2+x^2}&\text{si }x\leq-1\\[3pt]\dfrac{2x^2}{3-3x}&\text{si }x>-1\end{cases}\]\begin{tasks}(1)
  \task Estudia la continuidad de $f$ en $\mathbb R$.
  \task Calcula $\displaystyle\lim_{x\to-\infty}f(x)^{2x^2-1}$.
\end{tasks}

\ex{2} \pau{Islas Baleares, extraordinaria 2024} Indica el dominio de $g(x)=\dfrac{(x-1)\sqrt x}{x^2-1}$. ¿Es continua en su dominio? Si no lo es, ¿qué tipo de discontinuidad presenta?

\ex{2} \pau{Castilla y León, extraordinaria 2021} Estudia la continuidad de la función\[f(x)=\begin{cases}\dfrac{1-\cos x}{x}&\text{si }x\neq0\\[3pt]0&\text{si }x=0\end{cases}\]

\ex{2} \pau{Galicia, extraordinaria 2025} Estudia la continuidad en $x=0$ de la función\[f(x)=\begin{cases}x\,e^{4x}&\text{si }x<0\\[3pt]\dfrac{\ln(1+x)}{1+x}&\text{si }x\geq0\end{cases}\]

\ex{3} \pau{Navarra, ordinaria 2023} Estudia la continuidad en $\mathbb R$ de la función\[f(x)=\begin{cases}\dfrac{\cos^2(\pi x)-1}{1-x}&\text{si }x<1\\[3pt]\ln\left(x\,e^{x+1}\right)-2x&\text{si }x\geq1\end{cases}\]

\ex{2} \pau{Madrid, extraordinaria 2022} Estudia la continuidad en $x=0$ de $f(x)=\begin{cases}x&\text{si }x\leq0\\x\ln x&\text{si }x>0\end{cases}$.


% =====================================================
\seccio{7}{Funciones a trozos con parámetros}

\ex{2} Halla los parámetros para que las funciones sean continuas en todo su dominio:\begin{tasks}(2)
  \task $f(x)=\begin{cases}\dfrac{\sin(ax)}{x}&\text{si }x<0\\[3pt]x^2+3&\text{si }x\geq0\end{cases}$
  \task $g(x)=\begin{cases}\dfrac{e^{ax}-1}{x}&\text{si }x<0\\[3pt]\cos x+3&\text{si }x\geq0\end{cases}$
  \task $h(x)=\begin{cases}\dfrac{x^2+ax-10}{x-2}&\text{si }x<2\\[3pt]b+\ln(x-1)&\text{si }x\geq2\end{cases}$
  \task $k(x)=\begin{cases}\dfrac{\tan3x}{\sin x}&\text{si }-\frac\pi6<x<0\\[3pt]a&\text{si }x=0\\[3pt]\sqrt{x+b}&\text{si }x>0\end{cases}$
\end{tasks}\recorda{En cada punto de cambio igualamos los límites laterales y el valor de la función. Si un trozo da $\frac00$, resolvemos la indeterminación; si el límite tiene que ser finito y el denominador se anula, el numerador también se tiene que anular.}

\ex{2} Halla los parámetros para que las funciones sean continuas en $\mathbb R$:\begin{tasks}(1)
  \task $f(x)=\begin{cases}-2\sin x&\text{si }x\leq-\frac\pi2\\[3pt]a\sin x+b&\text{si }-\frac\pi2<x<\frac\pi2\\[3pt]\cos x&\text{si }x\geq\frac\pi2\end{cases}$
  \task $g(x)=\begin{cases}a\,e^{x}&\text{si }x\leq0\\[3pt]\dfrac{\ln(1+3x)}{x}&\text{si }0<x\leq1\\[3pt]b\sqrt{x+3}&\text{si }x>1\end{cases}$
\end{tasks}

\ex{2} Halla $k$ y $a$ para que las funciones sean continuas:\begin{tasks}(2)
  \task $f(x)=\begin{cases}\dfrac{\sqrt{x+1}-2}{x-3}&\text{si }x\neq3\\[3pt]k&\text{si }x=3\end{cases}$
  \task $g(x)=\begin{cases}x+e^2&\text{si }x\leq0\\[3pt](1+ax)^{1/x}&\text{si }x>0\end{cases}$
\end{tasks}

\ex{2} \pau{Andalucía, ordinaria 2022} Considera la función continua\[f(x)=\begin{cases}\dfrac1x&\text{si }x<-1\\[3pt]ax+b&\text{si }-1\leq x<1\\[3pt]\dfrac{x^2}{x+1}&\text{si }x\geq1\end{cases}\]\begin{tasks}(1)
  \task Calcula $a$ y $b$.
  \task Estudia y halla las asíntotas de la gráfica de $f$.
\end{tasks}

\ex{2} \pau{Aragón, extraordinaria 2024} Estudia la continuidad en $\mathbb R$, según los valores de $a$ y $b$, de\[f(x)=\begin{cases}a-\cos x&\text{si }x\leq0\\[3pt]x^2-b\sin\left(x+\frac\pi2\right)&\text{si }x>0\end{cases}\]

\ex{2} \pau{Aragón, ordinaria 2024} Estudia la continuidad en $\mathbb R$, según los valores de $a$, de\[f(x)=\begin{cases}\dfrac{e^{2x}-1}{x}&\text{si }x\neq0\\[3pt]a&\text{si }x=0\end{cases}\]

\ex{2} \pau{Aragón, ordinaria 2022} Determina $a$ para que $f$ sea continua en $\mathbb R$:\[f(x)=\begin{cases}\sqrt{-x}+e^3&\text{si }x\leq0\\[3pt](1-x)^{a/x}&\text{si }x>0\end{cases}\]

\ex{2} \pau{Aragón, extraordinaria 2021} Calcula los valores de $a\neq0$ para que $f$ sea continua:\[f(x)=\begin{cases}5-ax^2&\text{si }x\leq1\\[3pt]\dfrac{6}{ax}&\text{si }x>1\end{cases}\]

\ex{2} \pau{Islas Baleares, ordinaria 2022} Halla la condición que tienen que cumplir $a$ y $b$ para que $f$ sea continua:\[f(x)=\begin{cases}\dfrac{x^2+a}{2x-4}&\text{si }x\leq0\\[3pt]10x^2+x+b&\text{si }x>0\end{cases}\]

\ex{2} \pau{Islas Baleares, extraordinaria 2021} Considera $f(x)=\begin{cases}\dfrac{e^{ax}-1}{2x}&\text{si }x\neq0\\b&\text{si }x=0\end{cases}$.\begin{tasks}(1)
  \task Estudia la continuidad de $f$ en los puntos $x\neq0$.
  \task ¿Qué relación tiene que haber entre $a$ y $b$ para que $f$ sea continua en $x=0$?
\end{tasks}

\ex{2} \pau{Canarias, ordinaria 2021} Sea $f(x)=\dfrac{ax^2-2}{b-x}$.\begin{tasks}(1)
  \task Calcula $a$ y $b$ para que $f(-2)=2$ y $f$ sea continua en $\mathbb R-\{5\}$.
  \task Halla las asíntotas de $f$ si $a=-1$ y $b=-3$.
\end{tasks}

\ex{2} \pau{Cantabria, extraordinaria 2026} Considera\[f(x)=\begin{cases}x^2+3x+a^2&\text{si }x<1\\[3pt]\dfrac{\ln x}{x}+4a&\text{si }x\geq1\end{cases}\]\begin{tasks}(1)
  \task Calcula los valores de $a$ para los que $f$ es continua en todo su dominio.
  \task Calcula las asíntotas de $f$ en función de $a$.
\end{tasks}

\ex{2} \pau{Cantabria, ordinaria 2025} Considera $f(x)=\begin{cases}\dfrac{x^3-8}{x^2-4}&\text{si }x\neq2\\e^{k}&\text{si }x=2\end{cases}$.\begin{tasks}(1)
  \task Determina $k$ para que $f$ sea continua en $x=2$.
  \task Halla las asíntotas de $f$ y especifica su tipo.
\end{tasks}

\ex{2} \pau{Cantabria, extraordinaria 2020} Considera\[f(x)=\begin{cases}\sin x&\text{si }x\leq\frac\pi2\\[3pt]\dfrac2x+a&\text{si }x>\frac\pi2\end{cases}\]\begin{tasks}(1)
  \task Halla $a$ para que $f$ sea continua.
  \task Calcula $\displaystyle\lim_{x\to+\infty}f(x)$ y $\displaystyle\lim_{x\to-\infty}f(x)$.
\end{tasks}

\ex{2} \pau{Cantabria, extraordinaria 2018} Sea\[f(x)=\begin{cases}x+2&\text{si }x\leq-2\\[3pt]x^2+ax&\text{si }-2<x<0\\[3pt]2\sin x+b&\text{si }0\leq x\end{cases}\]\begin{tasks}(1)
  \task Determina $a$ y $b$ para que $f$ sea continua en $\mathbb R$.
  \task Si $a=3$ y $b=0$, clasifica la discontinuidad en $x=-2$.
\end{tasks}

\ex{2} \pau{Castilla-La Mancha, extraordinaria 2017} Determina $k$ para que $f$ sea continua en $x=0$:\[f(x)=\begin{cases}\left(\dfrac{x+1}{2x+1}\right)^{1/x}&\text{si }x<0\\[3pt]6x+k&\text{si }x\geq0\end{cases}\]

\ex{2} \pau{Castilla y León, extraordinaria 2024} ¿Qué valores tienen que tomar $a\in\mathbb R$ y $b\in\mathbb R-\{0\}$ para que $f$ sea continua en $\mathbb R$?\[f(x)=\begin{cases}\dfrac{\cos x-a}{bx^2}&\text{si }x\neq0\\[3pt]1&\text{si }x=0\end{cases}\]

\ex{2} \pau{Cataluña, ordinaria 2026} Determina los valores de $m$ que hacen que $f$ sea continua en todo su dominio:\[f(x)=\begin{cases}5e^{2x}&\text{si }x\leq0\\[3pt](x+m)^2+1&\text{si }0<x<2\\[3pt]1&\text{si }2\leq x\end{cases}\]

\ex{2} \pau{Extremadura, extraordinaria 2023} Halla $a$ y $b$ para que $f$ sea continua en $x=1$ y su gráfica pase por el punto $(-1,5)$:\[f(x)=\begin{cases}2x^2+ax+b&\text{si }x\leq1\\[3pt]\ln x&\text{si }x>1\end{cases}\]

\ex{2} \pau{Murcia, extraordinaria 2022} Considera $f(x)=\begin{cases}\dfrac{\ln x}{x-1}&\text{si }x>0,\ x\neq1\\a&\text{si }x=1\end{cases}$.\begin{tasks}(1)
  \task Calcula el límite de $f(x)$ cuando $x\to+\infty$.
  \task Determina $a$ para que $f$ sea continua en $x=1$.
\end{tasks}


% =====================================================
\seccio{8}{Teorema de Bolzano}

\ex{2} Demuestra que la ecuación $x^3+x-1=0$ tiene al menos una solución en el intervalo $(0,1)$ y hállala con una cifra decimal exacta.\recorda{\textbf{Bolzano}: si $f$ es continua en $[a,b]$ y $f(a)\cdot f(b)<0$, existe al menos un $c\in(a,b)$ con $f(c)=0$. Para afinar, parte el intervalo por la mitad y quédate con la mitad donde hay cambio de signo.}

\ex{2} Demuestra que:\begin{tasks}(1)
  \task la ecuación $x^5+x^3=1$ tiene alguna solución real, y halla un intervalo de longitud $0{,}1$ que la contenga;
  \task la ecuación $e^{-x}=x$ tiene alguna solución en $(0,1)$;
  \task las gráficas de $y=\cos x$ y $y=x$ se cortan en algún punto de $\left(0,\frac\pi2\right)$;
  \task la función $f(x)=x^3-x$ toma el valor $5$ en algún punto de $[1,2]$.
\end{tasks}

\ex{2} Razona si se puede aplicar el teorema de Bolzano y, si es el caso, qué podemos deducir:\begin{tasks}(1)
  \task $f(x)=\dfrac1x$ en $[-1,1]$
  \task $f(x)=\tan x$ en $\left[\frac\pi4,\frac{3\pi}4\right]$
  \task $f(x)=x^3+ax-1$ en $[0,1]$, con $a>0$
  \task $f(x)=x^2-1$ en $[-2,2]$
\end{tasks}

\ex{2} \pau{Islas Baleares, ordinaria 2025} La altitud (en metros) de un terreno en un tramo de 500 metros es $A(x)=-0{,}0001x^3+0{,}05x^2-4x+200$, con $x\in[0,500]$ la distancia recorrida horizontalmente. Demuestra que hay al menos un punto del tramo donde la altitud es 0.

\ex{2} \pau{Castilla-La Mancha, ordinaria 2023} Enuncia el teorema de Bolzano y utilízalo para justificar que $f(x)=x^3+6x^2+3x-10$ tiene al menos una raíz en $[0,2]$.

\ex{2} \pau{Castilla-La Mancha, extraordinaria 2017} Enuncia el teorema de Bolzano y comprueba si la ecuación $\cos x=2-x$ tiene alguna solución real en el intervalo $[0,2\pi]$.

\ex{2} \pau{Castilla y León, extraordinaria 2022} Averigua si la función $f(x)=x+\sin x-2$ se anula en algún punto del intervalo $\left[0,\frac\pi2\right]$.

\ex{2} \pau{Castilla y León, extraordinaria 2024} Dada la función $f(x)=e^x+x^3-2$, demuestra que $f$ se anula para algún valor de $x$ y que este valor es único.

\ex{2} \pau{Castilla y León, ordinaria 2017} Enuncia el teorema de Bolzano, interprétalo geométricamente y halla un intervalo en el que $P(x)=x^6+x^4-1$ tenga al menos una raíz.

\ex{2} \pau{Extremadura, ordinaria 2018} Enuncia el teorema de Bolzano y demuestra, usándolo, que la función $f(x)=x^3+x-3$ tiene una raíz real positiva.

\ex{2} \pau{Madrid, ordinaria 2022} Comprueba si $f(x)=\dfrac{x}{x^2+1}$ verifica las hipótesis del teorema de Bolzano en el intervalo $[-1,1]$.

\ex{3} \pau{Cataluña, ordinaria 2019} Considera la función $f(x)=\dfrac{2x^3-5x+4}{1-x}$.\begin{tasks}(1)
  \task Calcula el dominio y estudia la continuidad de $f$. ¿Tiene asíntota vertical?
  \task Observa que $f(-2)=-\frac23$, $f(0)=4$ y $f(2)=-10$. Razona si, a partir de esta información, podemos deducir que el intervalo $(-2,0)$ contiene un cero de la función. ¿Y el intervalo $(0,2)$? Halla un intervalo determinado por dos enteros consecutivos que contenga, como mínimo, un cero de la función.
\end{tasks}

\ex{2} \pau{Galicia, ordinaria 2026} Dibuja, si es posible, la gráfica de una función continua en $[a,b]$, con $f(a)f(b)<0$, y que corte el eje $X$ en más de un punto del intervalo $(a,b)$.

\ex{2} \pau{Navarra, ordinaria 2023} Se considera $f(x)=(x+1)\sin(\pi x)$. Comprueba que existe un valor $\alpha\in(0,1)$ tal que $f(\alpha)=\frac34$. Enuncia el resultado teórico utilizado.

\ex{3} \pau{Navarra, extraordinaria 2022} Sea $f(x)=\sin\left(\dfrac\pi4\ln\dfrac1x\right)$.\begin{tasks}(1)
  \task Demuestra que es continua en $\left[\frac1e,e\right]$.
  \task Demuestra que existe $\alpha\in\left(\frac1e,e\right)$ tal que $f(\alpha)=\dfrac{e\sqrt2}{1-e^2}$. Enuncia el resultado utilizado.
\end{tasks}

\ex{3} \pau{Navarra, extraordinaria 2021} Se considera $f(x)=\log_2\left[\sin\dfrac{\pi(x+1)}{4}+2^{\frac{x-5}{2}}\right]$.\begin{tasks}(1)
  \task Demuestra que es continua en $[6,7]$.
  \task Demuestra que existe $\alpha\in(6,7)$ tal que $f(\alpha)=0$.
\end{tasks}


% =====================================================
\seccio{9}{Asíntotas}

\ex{2} Calcula todas las asíntotas de las funciones siguientes y estudia los límites laterales en las verticales:\begin{tasks}(2)
  \task $f(x)=\dfrac{2x-1}{x+3}$
  \task $f(x)=\dfrac{2x^3-x}{x^2-1}$
  \task $f(x)=\dfrac{e^x}{x-1}$
  \task $f(x)=\sqrt{4x^2+x}$
  \task $f(x)=\ln\dfrac{x}{x-2}$
  \task $f(x)=x+\dfrac{\sin x}{x}$
\end{tasks}\recorda{AV: donde la función tiende a $\pm\infty$ (ceros de denominadores, extremos del dominio de un logaritmo\ldots). AH: $\lim_{x\to\pm\infty}f(x)=k$. AO $y=mx+n$: $m=\lim\frac{f(x)}{x}$, $n=\lim(f(x)-mx)$; estudia cada lado por separado.}

\ex{2} \pau{Andalucía, ordinaria 2021} La gráfica de $f(x)=\dfrac{ax^2+bx+2}{x-1}$ ($x\neq1$) tiene una asíntota oblicua que pasa por el punto $(1,1)$ y tiene pendiente 2. Calcula $a$ y $b$.

\ex{2} \pau{Cantabria, ordinaria 2026} Calcula los valores de $a$ para los que $f(x)=\dfrac{x^2+ax+1}{x+a^2}$ tiene como asíntota oblicua la recta $y=x-6$.

\ex{2} \pau{Cantabria, ordinaria 2026} Calcula, si las hay, las asíntotas de $f(x)=\dfrac{e^x-1}{e^{2x}-4}$.

\ex{2} \pau{Cantabria, extraordinaria 2018} ¿Para qué valor de $d$ tiene la función $\dfrac{x^d+1}{x-2}$ una asíntota oblicua en $+\infty$? Calcúlala.

\ex{2} \pau{Cantabria, extraordinaria 2020} Considera $f(x)=\dfrac{1-\cos x}{x}$. Calcula $\displaystyle\lim_{x\to0}f(x)$ y las asíntotas de $f$.

\ex{2} \pau{Canarias, extraordinaria 2020} Halla las asíntotas de la función $f(x)=\dfrac{x^3+5x^2}{x^2-1}$.

\ex{2} \pau{Castilla y León, extraordinaria 2022} Halla el dominio y las asíntotas de $f(x)=\dfrac{x^2}{2-x}$.

\ex{2} \pau{Islas Baleares, ordinaria 2025} Determina el dominio de $f(x)=\dfrac1x\ln(2x)$ y el comportamiento de la función en los extremos del dominio.

\ex{2} \pau{La Rioja, extraordinaria 2023} Halla el dominio y las asíntotas de $f(x)=\dfrac{1+4x^4-x^2}{x}$.

\ex{2} \pau{Madrid, extraordinaria 2022} Sea $f(x)=\begin{cases}\dfrac{2x+1}{x}&\text{si }x<0\\x^2-4x+3&\text{si }x\geq0\end{cases}$. Estudia la continuidad de $f$ y calcula, si las hay, las asíntotas horizontales y verticales.

\ex{2} \pau{País Vasco, extraordinaria 2025} Halla las asíntotas de $f(x)=\dfrac{x}{x^2-3x-4}$.

\ex{2} \pau{País Vasco, ordinaria 2024} Halla las asíntotas de $f(x)=\dfrac{x}{x^2-2x+1}$.

\ex{2} \pau{Extremadura, ordinaria 2018} Estudia el dominio y las asíntotas de $f(x)=\dfrac{1}{x^2-1}$.

\vfill
\begin{center}\color{muted}\small\faCheckCircle[regular]\ \ ¡Ánimo! Comprueba tus resultados con la calculadora de límites de los apuntes interactivos.\end{center}

\end{document}
